Para encerrar a batalha em $t_c=28$ dias com $E_0=21\,500$, $\lambda_F=0{,}0106$ e $\lambda_E=0{,}0544$ ($\alpha\approx0{,}02401$/dia):

Da condição $E(t_c)=0$,
\[
\tanh(\alpha t_c) = \frac{E_0}{F_0}\sqrt{\frac{\lambda_E}{\lambda_F}} \implies F_0 = \frac{E_0\,\sqrt{\lambda_E/\lambda_F}}{\tanh(\alpha t_c)}.
\]
\[
\tanh(0{,}02401\times28) = \tanh(0{,}6723) \approx 0{,}5866.
\]
\[
F_0 = \frac{21\,500\times2{,}265}{0{,}5866} \approx 83\,040\ \text{soldados}.
\]

As tropas restantes ao final seriam:
\[
F(t_c) = \sqrt{F_0^2 - \frac{\lambda_E}{\lambda_F}E_0^2} = \sqrt{83\,040^2 - 5{,}132\times21\,500^2} \approx \sqrt{4{,}523\times10^9} \approx 67\,250.
\]

As perdas dos EUA seriam $83\,040 - 67\,250 \approx 15\,790$ soldados, bem menores que as $\approx30\,680$ do cenário original (Problema~5.47, $F_0=54\,000$). Isso ilustra o efeito do quadrado de Lanchester: forças iniciais maiores reduzem as perdas relativas.